%% encoding: UTF-8
%% 首先读取模板的配置
%%%%%%%%%%%%%%%%%%%%%%%%%%%%%%%设置版面控制%%%%%%%%%%%%%%%%%%%%%%%%%%%%%%%%%%
\documentclass[a4paper, 10pt, notitlepage, oneside , twoside ]{article}
%\usepackage[colorlinks, linkcolor=red, anchorcolor=blue, citecolor=green, bookmarks]{hyperref}
%\hypersetup{xetex}
%\setlength{\parindent}{0em}
\usepackage[papersize={158.8mm,201.3mm},top=0cm,bottom=0.4cm,left=0cm,right=0cm]{geometry}
%\usepackage[papersize={200mm,255mm},top=3.37cm,bottom=2cm,left=2.06cm,right=2.06cm]{geometry}
%\usepackage[papersize={210mm,297mm},top=2.8cm,bottom=2.5cm,left=2.5cm,right=2.5cm]{geometry}                 % 设置版面，纸张空距
\usepackage{marvosym}
\usepackage{array}
\usepackage{makecell,multirow}
\usepackage{xcolor}
\usepackage{graphicx}
\RequirePackage{etoolbox,xpatch,pgfkeys,twoopt}
\usepackage{setspace}
%%%%%%%%%%%%%%%%%%%%%%%%%%%%%设置字体%%%%%%%%%%%%%%%%%%%%%%%%%%%%%
\usepackage[CJKchecksingle, CJKnumber]{xeCJK}                 % 用于字体
%\setCJKmainfont[BoldFont={FZLanTingHei-DB-GBK}]{Microsoft YaHei}
\setCJKmainfont[BoldFont={FZLanTingHei-DB-GBK}]{FZLanTingHei-L-GBK}
%\setmainfont{Microsoft YaHei}   % 设置中文/粗体/斜体/主要字体
\setmainfont{FZLanTingHei-L-GBK}
\punctstyle{hangmobanjiao}               % 设置行末半角
\newfontfamily\times{Times New Roman}
\newCJKfontfamily\kaishu{Adobe Kaiti Std}
\setCJKfamilyfont{zhyh}{Microsoft YaHei}
\newcommand*{\heiti}{\CJKfamily{zhyh}}
\setCJKfamilyfont{zhltxh}{FZLanTingHei-L-GBK}
\newcommand*{\lantxh}{\CJKfamily{zhltxh}}
\newfontfamily\arial{Arial}
\newcommand{\barial}{\bf\arial}
%\setmainfont[Mapping=tex-text]{Times New Roman}         % 设置衬线英文字体
\setsansfont[Mapping=tex-text]{FZLanTingHei-L-GBK}       % 设置无衬线英文字体
%\setmonofont{Courier New}                               % 设置等款英文字体
\newcommand{\chuhao}{\fontsize{42pt}{\baselineskip}\selectfont}
\newcommand{\xiaochuhao}{\fontsize{36pt}{\baselineskip}\selectfont}
\newcommand{\yihao}{\fontsize{28pt}{\baselineskip}\selectfont}
\newcommand{\erhao}{\fontsize{21pt}{\baselineskip}\selectfont}
\newcommand{\xiaoerhao}{\fontsize{18pt}{\baselineskip}\selectfont}
\newcommand{\sanhao}{\fontsize{15.75pt}{\baselineskip}\selectfont}
\newcommand{\sihao}{\fontsize{14pt}{\baselineskip}\selectfont}
\newcommand{\xiaosihao}{\fontsize{12pt}{\baselineskip}\selectfont}
\newcommand{\wuhao}{\fontsize{10.5pt}{\baselineskip}\selectfont}
\newcommand{\xiaowuhao}{\fontsize{9pt}{15pt}\selectfont}
\newcommand{\liuhao}{\fontsize{7.875pt}{\baselineskip}\selectfont}
\newcommand{\qihao}{\fontsize{7pt}{\baselineskip}\selectfont}
\newcommand{\bahao}{\fontsize{6.875pt}{\baselineskip}\selectfont}
\newcommand{\jiuhao}{\fontsize{6pt}{\baselineskip}\selectfont}
\newcommand{\RNum}[1]{\uppercase\expandafter{\romannumeral #1\relax}}
\usepackage{textgreek}

\definecolor{topcolor}{cmyk}{0.74,0.19,0.50,0}

%\RequirePackage{tcolorbox}
\usepackage[many]{tcolorbox}
\tcbuselibrary{skins}
\newenvironment{LRaside}[2][0.7]{%
  \newcommand{\asidebreak}{\tcblower}
  \begin{tcolorbox}[title={#2},
  righthand width=#1\textwidth, sidebyside, enhanced, colback=white,
  colbacktitle=white,coltitle=gray2,coltext=gray2,colframe=topcolor,arc=0pt,
  segmentation hidden, boxrule=1pt, titlerule=-1pt, toptitle=\bigskipamount, top=\bigskipamount, bottom=\bigskipamount,
  segmentation empty,
  fonttitle=\sffamily\large,fontupper=\sffamily,fontlower=\sffamily,parbox=false,
  before skip=\bigskipamount,after skip=\bigskipamount,  attach boxed title to top left={yshift=-3mm,xshift=5mm},
  boxed title style={colframe=white,arc=0pt,outer arc=0pt,boxrule=0pt,toprule=0pt,left=-1pt,right=-1pt},]
  }{%
%\parindent0pt}{%
  \end{tcolorbox}
}

\newenvironment{CRaside}[2][0.7]{%
  \newcommand{\asidebreak}{\tcblower}
  \begin{tcolorbox}[title={#2},
  righthand width=#1\textwidth, sidebyside, enhanced, colback=white,
  colbacktitle=white,coltitle=white,coltext=gray2,colframe=topcolor,arc=0pt,
  segmentation hidden, boxrule=1pt, titlerule=-1pt, toptitle=\bigskipamount, top=\bigskipamount, bottom=\bigskipamount,
  segmentation empty,
  fonttitle=\sffamily\large,fontupper=\sffamily,fontlower=\sffamily,parbox=false,
  before skip=\bigskipamount,after skip=\bigskipamount,  attach boxed title to top center={yshift=-3mm,xshift=5mm},
  boxed title style={colframe=topcolor,colback=topcolor,arc=0pt,outer arc=0pt,boxrule=0pt,toprule=0pt,left=12pt,right=12pt},]
  }{%
%\parindent0pt}{%
  \end{tcolorbox}
}

%%%%%%%%%%%%%%%%%%%%%%%%%%%%%设置图像%%%%%%%%%%%%%%%%%%%%%%%%%%%%%
                   % 用于图像编号
\graphicspath{{_drugRisk/}{_logo/}{_phenoDegree/}{disease_pic/}{drug_pic/}{pheno_pic/}{cores/}}
\usepackage{tikz}
\usetikzlibrary{calc}
\usetikzlibrary{shapes.geometric, arrows}
\usetikzlibrary{circuits.logic.US,circuits.logic.IEC,circuits.ee.IEC}
\usepackage{verbatim}
\usepackage{subfig}
%\usepackage{mathpazo}
\usepackage{wrapfig}
\usepackage{pgfplots}
\pgfplotsset{compat=newest}% <-- moves axis labels near ticklabels (respects tick label widths)
\usepackage{wallpaper}%%%%%%%%%%%%%%%%背景%%%%%%%%%%%%%%%%%%
%%%%%%%%%%%%%%%%%%%%%%%%%%mdframed%%%%%%%%%%%%%%%%%%%%

\usepackage{fancyhdr}
\pagestyle{fancy}
\renewcommand{\headrulewidth}{0pt}
\renewcommand{\footrulewidth}{0pt}
%%%%%%%%%%%%%%%页码修改[RO]居右[CO]居中[LO]居左%%%%%%%%%%%

%\cfoot{}
%\fancyfoot[LO,RE]{}
%\fancyhead[LO,LE]{\raisebox{-1.5ex}{\includegraphics{logo}}}

\usepackage{paralist}
\usepackage{float,longtable,hhline}

%\definecolor{topcolor}{cmyk}{0.28,0,1,0}
%\definecolor{topcolor}{cmyk}{.20,0,1,0}
\definecolor{green}{cmyk}{0.71,0.15,0.89,0}
%\definecolor{green2}{RGB}{196,224,177}
\definecolor{green2}{cmyk}{0.28,0,0.39,0}
\definecolor{green3}{cmyk}{0.57,0,0.75,0}
\definecolor{green4}{cmyk}{0.77,0.23,1,0}
\definecolor{yellow}{RGB}{238,196,133}
%\definecolor{yellow1}{RGB}{243,225,0}
\definecolor{yellow1}{cmyk}{0.00,0.50,1.00,0.00}
\definecolor{yellow2}{RGB}{243,152,0}
\definecolor{columncolor}{cmyk}{0,0,0,.08}%{RGB}{247,248,248}
\definecolor{red}{cmyk}{0.07,0.79,0.66,0}
\definecolor{red2}{cmyk}{0.1,0.58,0.47,0}
\definecolor{orange}{RGB}{245,130,31}
\definecolor{orange2}{cmyk}{0.07,0.27,0.51,0}
%\definecolor{gray1}{RGB}{210,210,210}
%\definecolor{gray2}{RGB}{144,144,144}
\definecolor{gray3}{RGB}{100,100,100}
\definecolor{gray0}{RGB}{244,244,244}
\definecolor{blue1}{RGB}{65,100,140}
\definecolor{blue2}{RGB}{50,168,225}
\definecolor{gray2}{cmyk}{0.00,0.00,0.00,0.70}
\definecolor{gray1}{cmyk}{0.00,0.00,0.00,0.30}
\definecolor{black100}{cmyk}{0.00,0.00,0.00,1.00}
%\arrayrulecolor{gray2}

\def\blueC{{\color{blue}C}}
\def\blueG{{\color{blue}G}}
\def\redA{{\color{red}A}}
\def\redT{{\color{red}T}}
\def\yellowA{{\color{yellow}A}}
\def\yellowT{{\color{yellow}T}}
\def\redrule{\raisebox{.5ex}{~~~\hspace*{-3cm}{\color{red}\rule{2.3cm}{1.5pt}}}~}

\def\sun{\tikz\draw[yellow1,fill=yellow1](0,0) circle[radius=4.5pt];
\hspace*{-0.5cm}\tikz\draw[yellow1,dotted](0,0) circle[radius=5pt];}


\def\trianglea{\tikz \draw[scale=0.07,color=yellow1,fill=yellow1] (1,0) -- (-1,0) -- (0,1.732) -- cycle;}
\def\triangleb{\tikz \draw[scale=0.07,color=yellow1,fill=yellow1] (1,0) -- (-1,0) -- (0,-1.732) -- cycle;}
\usepackage{tabularx}

%定义循环插图

\newcount\tmp

\newcommand{\pdfchatu}[2]{% need this to prevent extra vertical space
    % #1: image path and core part of name
    % #2: maximum number
    \tmp=0
    \loop
        % increment dummy counter
        \advance\tmp by 1
        % include the image
        \includegraphics[width=.35cm]{xiaoren00.pdf}
        % repeat the loop provided the counter is within specified bound
        \ifnum\tmp<#2
    \repeat
}


\newcount\tmp

\newcommand{\pdfchatuoliver}[2]{% need this to prevent extra vertical space
    % #1: image path and core part of name
    % #2: maximum number
    \tmp=0
    \loop
        % increment dummy counter
        \advance\tmp by 1
        % include the image
        \includegraphics[width=.35cm]{xiaoren010.pdf}
        % repeat the loop provided the counter is within specified bound
        \ifnum\tmp<#2
    \repeat
}
\newcount\tmp

\newcommand{\pdfchatuorange}[2]{% need this to prevent extra vertical space
    % #1: image path and core part of name
    % #2: maximum number
    \tmp=0
    \loop
        % increment dummy counter
        \advance\tmp by 1
        % include the image
        \includegraphics[width=.35cm]{xiaoren110.pdf}
        % repeat the loop provided the counter is within specified bound
        \ifnum\tmp<#2
    \repeat
}

\newtcolorbox{tctabularx}[1]{fonttitle=\bfseries, colback=white, colframe=topcolor,#1,sharp corners,before upper app={\rowcolor{topcolor}}}    % 表头和minipage左上角设置为直角

\makeatletter
\tcbset
    {
tabularx*/.style 2 args={%
    boxsep=0pt,top=0pt,bottom=0pt,leftupper=0pt,rightupper=0pt,
    toptitle=1mm,bottomtitle=1mm,boxrule=0.5mm,
    before upper={\arrayrulecolor{black100}\def\arraystretch{1}#1%
      \tcb@hack@currenvir\tabularx{\linewidth}{#2}},
    after upper=\endtabularx\arrayrulecolor{gray}}
    }
\makeatother


%\arrayrulecolor{gray2}
